

Le schéma de base de données suivant sera utilisé pour l'ensemble du
sujet. Il permet de gérer une plateforme de vente de billets de spectacles en ligne :

\begin{itemize}
  \item Personne(id, nom, prénom, age)
  \item Salle(id, nom, adresse, nbPlaces)
  \item Spectacle (id, idArtiste, idSalle, dateSpectacle, nbBilletsVendus)
   \item Billet (id, idSpectateur, idSpectacle, catégorie, prix)
\end{itemize}

Notez que les artistes et les spectateurs sont des personnes. Tous les attributs \textbf{id}
sont des clés primaires, ceux commençant par \textbf{id} sont des clés étrangères.

\section{Conception (5 points)}

\begin{enumerate}
    \item Donnez le modèle entité / association correspondant au schéma de notre base.
     \item Pensez-vous qu'une association a été réifiée? Si oui, laquelle?
     \item Donnez la commande de création de la table \texttt{Spectacle}
     \item On veut qu'un même spectateur puisse avoir au plus un billet pour un spectacle. 
        Quelle dépendance fonctionnelle cela ajoute-t-il? Que proposez-vous pour en tenir compte dans le schéma\,?
     
\end{enumerate}



\begin{solution}

\begin{minted}{sql}
SELECT p.nom
FROM Personne as p, SPECTACLE s, SALLE l
WHERE p.id = s.idArtiste AND s.idSalle = l.id 
AND l.nom = 'Bercy' AND p.age >= 70
\end{minted}
\end{solution}
\section{SQL  (6 points)}

\textbf{Question de cours}. Deux requêtes syntaxiquement différentes sont équivalentes si
\begin{itemize}
    \item Elles peuvent donner le même résultat
   \item  Elles donnent toujours le même résultat
\end{itemize}

   Indiquez la bonne réponse.
   

   Donnez les expressions SQL des requêtes suivantes
   
\begin{enumerate}
   
  \item Noms des artistes âgé d'au moins 70 ans qui ont donné au moins un spectacle dans la salle 'Bercy' 

   \item Donnez le nom des artistes, des salles où ils se sont produits, triés sur le nombre de billets vendus

  \item Quelles personnes n'ont jamais acheté de billet\,?
  
\item  Donnez la liste des spectateurs ayant acheté au moins 10 billets,
          affichez le nom, le nombre de billets et le prix total.
 
\end{enumerate}


\begin{solution}

\begin{minted}{sql}
SELECT p.nom
FROM Personne as p, SPECTACLE s, SALLE l
WHERE p.id = s.idArtiste AND s.idSalle = l.id 
AND l.nom = 'Bercy' AND p.age >= 70
\end{minted}

\begin{minted}{sql}
SELECT p.nom as nomArtiste, s.nom as nomSalle, s.nbBilletsVendus
FROM Peronne as p, SPECTACLE as sp, Salle as s
WHERE p.id = sp.idArtiste
and sp.idSalle = s.id
order by s.nbBilletsVendus
\end{minted}

\begin{minted}{sql}
SELECT prénom, nom
FROM Personne as p
WHERE not exists (select '' from Billet where idSpectateur = p.id)

\end{minted}

\begin{minted}{sql}
SELECT p.nom, count(*), sum(prix)
FROM Personne as p, Billet as b
WHERE p.id = b.idSpectateur
GROUP BY p.id, p.nom
HAVING count(*) >= 10
\end{minted}
\end{solution}

\section{Algèbre (5 points)}

Voici 3 relations.

\medskip

\begin{minipage}[b]{5cm}
\begin{tabular}{llll}
A  & B   & C   & D\\ \hline
1  &  0  &  1   & 2\\
4   & 1  &  2   & 2\\
6    &2  &  6   & 3\\
7   & 1   & 1   & 3\\
1   & 0  &  1   & 0\\
1   & 1 &   1     &0\\ \hline
\end{tabular}

Relation R
\end{minipage}\
\begin{minipage}[b]{5cm}
\begin{tabular}{lll}
 A  &B  & E\\ \hline
1  & 0  & 2\\
4  & 1  & 2\\
1 &  0 &  7\\
8  & 6  & 6\\ \hline
\end{tabular}

Relation S
\end{minipage}
\
\begin{minipage}[b]{5cm}
\begin{tabular}{lll}
 A  &B  & E\\ \hline
1  & 0  & 2\\
4  & 2  & 2\\
1 &  0 &  7\\
8  & 6  & 5\\
8  & 5  & 6\\ \hline
\end{tabular}

Relation T
\end{minipage}


\medskip

\begin{enumerate}
  \item Pour quelle requête le résultat contient-il plus d’un nuplet\,?  Attention: souvenez-vous que l'opérateur
    de projection élimine les doublons.

    \begin{enumerate}
        \item  $\pi_{A,C} (\sigma_{B=0}(R))$
        \item $\pi_{A,C} (\sigma_{D=0}(R))$
       \item  $\pi_{A,C} (\sigma_{B=0}(R) \cup \sigma_{D=0}(R))$
    \item   $\pi_{A,C} (\sigma_{A=C}(R)$
    \end{enumerate}
    
    \item Combien de nuplets retourne la requête $\pi_{A,B,E} (S \Join_{A=A \land B=B} R) $?
    
           \begin{enumerate}
        \item  2
        \item 3
        \item 4
        \item 5
    \end{enumerate}
   \item  Combien de nuplets retourne la requête $R \Join_{A=A \land B=B} (S \cup T)$ ?
       \begin{enumerate}
        \item  3
        \item 5
        \item 7
        \item 9
    \end{enumerate}
   
  \item Donnez l'expression algébrique pour la requête ``Noms et prénoms des spectateurs qui ont 
      acheté au moins une fois un billet de plus de 500 euros.''
  \item Donnez l'expression algébrique pour la requête ``Id des personnes qui ne sont pas des artistes''
  
  $$pi_id (Personne) - rho_{idArtiste -> id} (pi_idArtiste (Spectacle))$$
\end{enumerate}



\begin{solution}
    $$\pi_{nom, prenom}(Personne \underset{id=idSpectateur}{\bowtie} \sigma_{prix > 500}(BILLET))$$

   $$\pi_{id}(Personne - \pi_{idArtiste} (Billet))$$
\
\end{solution}

\section{Transactions (4 points)}

On prend le schéma donné au début de l'examen.
      
\begin{enumerate}

  \item \textbf{Question de cours}. Quel est le principal inconvénient du mode \texttt{serializable}\,?

  \begin{itemize}
    \item Le système ne peut exécuter qu’une transaction à la fois
    \item Certaines transactions sont mises en attente avant de pouvoir finir
    \item Certaines transactions sont rejetées par le système
\end{itemize}

   Indiquez la bonne réponse
   
\item Exprimez au moins un critère de cohérence de cette base 
 

\begin{solution}
Réponse: le nombre de billets vendus (\texttt{nbBilletsVendus} dans \texttt{Spectacle}) pour un spectacle est égal au nombre
de lignes dans \texttt{Billet} pour ce même spectacle.

Autre possibilité\,: pas plus de billets vendus que de places dans la salle.
\end{solution}

\item On écrit une procédure de réservation d'un billet pour un spectacle.

\begin{verbatim}
Reserver (idSpectateur, idSpectacle, categorie, prix)
   /// Requête SQL de mise à jour du spectacle
   requete A
   // Requête SQL d'insertion du billet
   requete  B
\end{verbatim}

Donnez les deux requêtes SQL, indiquez l'emplacement des \texttt{commit} et \texttt{rollback} 
et justifiez-les brièvement.

\begin{solution}
Réponse:  un \texttt{update} de \texttt{Spectacle} pour incrémenter \texttt{nbBilletsVendus},
et un \texttt{insert} dans \texttt{Billet}. Le \texttt{commit} vient à la fin car c'est seulement à ce moment
que la base est à nouveau dans un état cohérent. 
\end{solution}

\end{enumerate}

